\documentclass[11pt]{article}
\usepackage[margin=1in]{geometry}
\usepackage{amsmath,amssymb,amsthm}
\usepackage[hidelinks]{hyperref}
\newtheorem{theorem}{Theorem}
\newtheorem{lemma}{Lemma}
\DeclareMathOperator{\Lip}{Lip}
\title{A unique upper transition attractor for bounded similarity families}
\author{Principia Math in Progress\\MathDB 344775}
\date{12 September 2026}
\begin{document}
\maketitle
\begin{abstract}
We prove Vince's Conjecture 8.1 for every bounded one-parameter family in
the source's sense, in every finite dimension. The powers of the distinguished
isometry need not be periodic. Damped powers converge as a compact set of
operators, and regrouping the original maps gives a contraction constant
bounded strictly below one throughout the limiting process.
\end{abstract}

\paragraph{Task 1. Prove uniqueness of the upper transition attractor.}
Here ``bounded'' has the technical meaning of Definition 8.1 in \cite{V}:
the family consists of similarities, each map has a fixed point independent
of the parameter, and exactly one underlying similarity has the largest
scaling ratio. It is not an assumption merely about bounded sets for an
arbitrary affine family.

By rescaling the parameter so that the threshold is $1$ and translating
the special fixed point to $0$, the source's Proposition 5.4 gives the form
\begin{equation}\label{normal}
 \mathcal F_t=\{g_t,f_{1,t},\ldots,f_{m,t}\},\quad
 g_t(x)=tUx,\quad f_{i,t}(x)=tL_i(x-q_i)+q_i,
\end{equation}
where $U\in O(d)$, $q_i\in\mathbb R^d$, and
$\rho:=\max_i\|L_i\|<1$. For the original problem each $L_i$ is a similarity;
the proof below only uses its displayed norm bound. Let $A_t$ be the unique
attractor for $0\le t<1$, with distances between compact sets measured by
the Hausdorff metric $h$.

\begin{theorem}\label{main}
For the family \eqref{normal}, $A_t$ has a limit $A_*$ as $t\uparrow1$.
Consequently it has a unique upper transition attractor, and the map
$t\mapsto A_t$ is uniformly continuous on $(0,1)$.
\end{theorem}

We first isolate the compactness step. Operator norms are Euclidean norms,
and $h_{\rm op}$ is the resulting Hausdorff metric on compact sets of matrices.

\begin{lemma}\label{powers}
For $U\in O(d)$ set
\[
 H=\overline{\{U^n:n\ge0\}},\quad
 E_t=\{0\}\cup\{t^nU^n:n\ge0\}\quad(0<t<1),\quad
 D=\{rQ:0\le r\le1,\ Q\in H\}.
\]
Then $H$ is a compact group, $E_t$ and $D$ are compact, and
$h_{\rm op}(E_t,D)\to0$ as $t\uparrow1$.
\end{lemma}
\begin{proof}
The set $H$ is a closed subsemigroup of the compact group $O(d)$.
Choose an increasing sequence $n_k$ such that $U^{n_k}$ converges. Then
$U^{n_{k+1}-n_k}\to I$, so $U^{n_{k+1}-n_k-1}\to U^{-1}$, and $U^{-1}\in H$.
Thus all integral powers belong to $H$, and continuity of inversion shows
that $H$ is a group. Compactness of $D$ follows from its parametrization
by $[0,1]\times H$. For fixed $t<1$, the only possible accumulation point
of $t^nU^n$ as $n\to\infty$ is $0$, so $E_t$ is compact.

Fix $\varepsilon>0$. Density of the nonnegative powers and compactness give
an integer $L\ge0$ such that $I,U,\ldots,U^L$ is an $\varepsilon$-net of $H$.
For every $N\ge0$, the list $U^N,\ldots,U^{N+L}$ is also an
$\varepsilon$-net, since left multiplication by $U^N$ is an isometry of $H$.

Given $rQ\in D$ with $r>0$, choose
$N=\lfloor\log r/\log t\rfloor$, so $t^{N+1}<r\le t^N$.
Some $j\in\{0,\ldots,L\}$ satisfies $\|U^{N+j}-Q\|\le\varepsilon$.
It follows that
\begin{align*}
 \|t^{N+j}U^{N+j}-rQ\|
 &\le \varepsilon+|t^{N+j}-r|\\
 &\le \varepsilon+(1-t^L)+(1-t).
\end{align*}
The case $r=0$ is exact, since $0\in E_t$. Conversely $E_t\subseteq D$.
Thus
\begin{equation}\label{netbound}
 h_{\rm op}(E_t,D)\le\varepsilon+(1-t^L)+(1-t).
\end{equation}
First let $t\uparrow1$ with $\varepsilon,L$ fixed, then let $\varepsilon$
decrease to zero. No periodicity assumption has been used.
\end{proof}

\begin{proof}[Proof of Theorem~\ref{main}]
Choose $C\ge\sup_{i,\,0\le t\le1}\|f_{i,t}(0)\|$ and a radius
$R>\max\{1,C/(1-\rho)\}$. The closed ball $B=\overline B(0,R)$ is
invariant under all the maps in \eqref{normal}, including at $t=1$:
\[
 \|f_{i,t}(x)\|\le\rho R+C\le R,\qquad \|g_t(x)\|\le R
 \quad(x\in B).
\]
In particular $A_t\subseteq B$. Define compact families of maps on $B$ by
\[
 \mathcal M_t=\{x\mapsto M f_{i,t}(x):M\in E_t,\ 1\le i\le m\},\qquad
 \mathcal M_* =\{x\mapsto M f_{i,1}(x):M\in D,\ 1\le i\le m\}.
\]
Each family includes the constant-zero map. Every map in these families
has Lipschitz constant at most $\rho$, uniformly in $t$ and $M$.
Their Hutchinson operators $\Phi_t,\Phi_*$ on the complete metric space
$\mathcal K(B)$ of nonempty compact subsets of $B$ are therefore
$\rho$-contractions. Their values are compact because the parameter
sets and the input sets are compact. The contraction theorem gives a unique
fixed compact set for each operator; denote the fixed set of $\Phi_*$ by $A_*$.

We claim that the fixed set of $\Phi_t$ is exactly $A_t$. Put
$F_t(K)=\bigcup_i f_{i,t}(K)$. Iterating the original invariance equation gives
\begin{equation}\label{unfold}
 A_t=\bigcup_{n=0}^{k-1}g_t^n(F_t(A_t))\ \cup\ g_t^k(A_t)
 \qquad(k\ge1).
\end{equation}
Also $0\in A_t$, since repeated application of $g_t$ to any point of $A_t$
converges to $0$. Every set in the union on the right lies in $A_t$.
Conversely, if a point of $A_t$ lies in none of these sets for any $n$,
then \eqref{unfold} places it in $g_t^k(A_t)$ for every $k$, so its norm is
at most $t^kR$ for every $k$ and it is zero. Since $g_t^n(x)=t^nU^nx$,
\[
 A_t=\{0\}\cup\bigcup_{n\ge0,\,i}g_t^n(f_{i,t}(A_t))=\Phi_t(A_t).
\]
This proves the claim, including the required infinite tails.

Let
\[
 \eta_t=h_{\rm op}(E_t,D),\qquad
 e_t=\max_i\sup_{x\in B}\|f_{i,t}(x)-f_{i,1}(x)\|.
\]
Lemma~\ref{powers} gives $\eta_t\to0$. Formula \eqref{normal} gives
$e_t\le(1-t)\max_i\sup_{x\in B}\|L_i(x-q_i)\|\to0$.
Pairing close matrices in the two coefficient sets, using
$f_{i,1}(B)\subseteq B$ and $\|M\|\le1$, shows that for every
$K\in\mathcal K(B)$,
\[
 h(\Phi_t(K),\Phi_*(K))\le e_t+R\eta_t.
\]
The fixed-point equations now imply
\[
 h(A_t,A_*)\le \rho h(A_t,A_*)+e_t+R\eta_t,
 \qquad
 \boxed{\ h(A_t,A_*)\le\frac{e_t+R\eta_t}{1-\rho}\ \longrightarrow\ 0.\ }
\]
Hence the upper transition attractor is unique.

For completeness, on every interval $[0,b]$ with $b<1$, the original
Hutchinson operators have a common contraction factor at most $b$ and depend
uniformly continuously on $t$ on $B$. The same fixed-point estimate proves
continuity of $A_t$ there, including at $0$, where the attractor is the
finite set of constant-map values. Together with the established limit at
$1$, this extends $t\mapsto A_t$ continuously to $[0,1]$ and thus proves
uniform continuity. Rescaling back gives the assertion on $(0,t_0)$.
\end{proof}

If the original family has only the distinguished map, its attractor is
the special fixed point for every $t<1$, so the conclusion holds immediately.
The limiting set above need not be an attractor of the original threshold
IFS $\mathcal F_1$; only uniqueness of the upper transition attractor is asserted.

The final 2022 article \cite{LSSV}, p.~5418, explicitly describes its
periodic-isometry theorem as a partial answer to Vince's conjecture. Its
Question 6.2 on p.~5425 leaves the nonperiodic case open. The compact
operator-set limit in Lemma~\ref{powers} supplies that missing step for
the finite-dimensional source problem.

\paragraph{Task 2. Quantify convergence for nonperiodic rotations.}
Use covering times of $\{U^n\}$ in $H$ and \eqref{netbound} to derive
rates for $h(A_t,A_*)$. Determine the sharp rates under Diophantine
conditions on the rotation angles and contrast them with the linear bound
available when $U$ is periodic.

\paragraph{Task 3. Determine useful infinite-dimensional hypotheses.}
Extend the argument to Banach spaces under compactness assumptions on the
isometry powers. Determine whether compactness of the relevant point
orbits can replace compactness in operator norm.

\paragraph{Task 4. Treat several maps at the critical scaling ratio.}
Determine when multiple distinguished similarities admit a unique limiting
attractor. Analyze the roles of a common fixed point, the group they
generate, and boundedness of the attractors near the threshold.

\paragraph{Task 5. Analyze the geometry of the limiting set.}
Use the explicit family $\mathcal M_*$ to characterize the radial,
symmetry and interior properties of $A_*$. Determine when this set agrees
with the lower transition attractor and when the containment is strict.

\begin{thebibliography}{9}
\bibitem{V}
A. Vince, \emph{Thresholds for one-parameter families of affine iterated
function systems}, Nonlinearity \textbf{33} (2020), 6541--6563.
\href{https://people.clas.ufl.edu/avince/files/Vince_2020_Nonlinearity.pdf}{Final journal PDF};
\href{https://arxiv.org/abs/1912.03220}{arXiv:1912.03220}.
Definition 8.1 and Conjecture 8.1 appear on pp.~6557--6558.
\bibitem{LSSV}
K. Le\'{s}niak, N. Snigireva, F. Strobin and A. Vince,
\emph{Transition phenomena for the attractor of an iterated function system},
Nonlinearity \textbf{35} (2022), 5396--5426.
\href{https://doi.org/10.1088/1361-6544/ac8af1}{doi:10.1088/1361-6544/ac8af1}.
Theorem 5.2, p.~5418, and Question 6.2, p.~5425.
\end{thebibliography}
\end{document}
